\documentclass[11pt]{amsart}
\usepackage{amsmath,amssymb,amsthm}

\newtheorem{theorem}{Theorem}[section]
\newtheorem{lemma}[theorem]{Lemma}
\newtheorem{proposition}[theorem]{Proposition}
\newtheorem{corollary}[theorem]{Corollary}
\newtheorem{conjecture}[theorem]{Conjecture}
\theoremstyle{definition}
\newtheorem{definition}[theorem]{Definition}
\theoremstyle{remark}
\newtheorem*{remark}{Remark}

\title{Repeated degrees in finite simple graphs}
\author{M. K. Sample}
\address{Institute for Hypothetical Combinatorics (fictional)}
\email{m.k.sample@example.org}
\date{}

\begin{document}

\begin{abstract}
It is a well-known party puzzle that among any two or more people, two have
the same number of acquaintances present. We give a careful proof of this fact
for finite simple graphs, based on the observation that a graph cannot contain
both an isolated and a dominating vertex. We also collect two elementary
sufficient conditions for connectivity and comment on the characterisation of
degree sequences due to Erd\H{o}s and Gallai.
\end{abstract}

\maketitle

\section{Introduction}

The pigeonhole argument showing that no finite simple graph on at least two
vertices has all degrees distinct appears in many textbooks, e.g.~\cite{Bol}.
Our goal is to present it (Theorem~\ref{thm:main}) in a form where every
step is isolated as a separate statement, together with some related facts on
degrees and connectivity that are frequently used alongside it.

\section{Degrees}

\begin{definition}\label{def:graph}
A \emph{graph} $G=(V,E)$ consists of a finite set $V$ of \emph{vertices} and a
set $E$ of $2$-element subsets of $V$, the \emph{edges}. We write $n=|V|$. For
$v\in V$, the \emph{neighbourhood} of $v$ is $N(v)=\{u\in V:\{u,v\}\in E\}$ and
the \emph{degree} of $v$ is $\deg v=|N(v)|$. A vertex is \emph{isolated} if its
degree is $0$ and \emph{dominating} if its degree is $n-1$. The graph is
\emph{connected} if any two vertices are joined by a path.
\end{definition}

\begin{lemma}[Handshake lemma]\label{lem:hand}
For every graph $G$, $\sum_{v\in V}\deg v=2|E|$.
\end{lemma}

\begin{proof}
Count the pairs $(v,e)$ with $v\in e\in E$: each edge contains exactly two
vertices, and each vertex $v$ lies in exactly $\deg v$ edges.
\end{proof}

\begin{corollary}\label{cor:odd}
In every graph the number of vertices of odd degree is even.
\end{corollary}

\begin{proof}
Reduce Lemma~\ref{lem:hand} modulo $2$: the sum of the odd degrees is congruent
to the number of vertices of odd degree, and it is even.
\end{proof}

\begin{lemma}\label{lem:range}
For every vertex $v$ of a graph on $n$ vertices, $0\le \deg v\le n-1$.
\end{lemma}

\begin{proof}
$N(v)\subseteq V\setminus\{v\}$, since edges have two distinct elements.
\end{proof}

\begin{lemma}\label{lem:notboth}
No graph has both an isolated vertex and a dominating vertex.
\end{lemma}

\begin{proof}
Suppose $u$ is dominating and $w$ is isolated. Then $u\ne w$, because their
degrees $n-1$ and $0$ differ. As $u$ is adjacent to every other vertex, $w\in
N(u)$, so $u\in N(w)$ and $\deg w\ge 1$, a contradiction.
\end{proof}

\begin{theorem}[Main theorem]\label{thm:main}
Let $G$ be a graph with $n\ge 2$ vertices. Then there are two distinct vertices
$u\ne w$ with $\deg u=\deg w$.
\end{theorem}

\begin{proof}
By Lemma~\ref{lem:range} every degree lies in $\{0,1,\dots,n-1\}$, a set of $n$
values. By Lemma~\ref{lem:notboth} the values $0$ and $n-1$ do not both occur,
so the degrees of $G$ take at most $n-2$ distinct values. Since there are $n$
vertices, the pigeonhole principle gives two distinct vertices with the same
degree.
\end{proof}

\begin{proposition}\label{prop:distinct}
Let $G$ be a graph on $n\ge 2$ vertices whose degrees are pairwise distinct.
Then $G$ has exactly one isolated vertex and exactly one dominating vertex.
\end{proposition}

\begin{proof}
The $n$ distinct degrees lie in the $n$-element set $\{0,\dots,n-1\}$ by
Lemma~\ref{lem:range}, so each value is taken exactly once; in particular $0$
and $n-1$ are each taken by exactly one vertex.
\end{proof}

\section{Connectivity}

\begin{proposition}\label{prop:mindeg}
Let $G$ be a graph on $n\ge 1$ vertices in which every vertex has degree at
least $\frac{n-2}{2}$. Then $G$ is connected.
\end{proposition}

\begin{proof}
Let $u\ne w$ be non-adjacent vertices. Then $N(u)$ and $N(w)$ are subsets of
$V\setminus\{u,w\}$, which has $n-2$ elements, while
$|N(u)|+|N(w)|\ge \frac{n-1}{2}+\frac{n-1}{2}=n-1>n-2$. Hence $u$ and $w$ have a
common neighbour, so any two vertices are joined by a path of length at most
$2$.
\end{proof}

\begin{proposition}\label{prop:edges}
Let $G$ be a graph on $n\ge 1$ vertices with $|E|\ge\binom{n-1}{2}$. Then $G$
is connected.
\end{proposition}

\begin{proof}
If $G$ is not connected, split $V=A\sqcup B$ with $A,B$ nonempty and no edge
between $A$ and $B$. With $a=|A|$ we get
$|E|\le\binom{a}{2}+\binom{n-a}{2}\le\binom{n-1}{2}$, since the middle
expression is a convex function of $a\in[1,n-1]$, maximised at the endpoints.
\end{proof}

\section{Degree sequences}

The \emph{degree sequence} of $G$ is the list of its degrees in non-increasing
order. We recall the classical characterisation of degree sequences.

\begin{theorem}[Erd\H{o}s--Gallai~\cite{EG}]\label{thm:EG}
A sequence $d_1\ge d_2\ge\dots\ge d_n\ge 0$ of integers is the degree sequence
of a graph if and only if $\sum_{i=1}^n d_i$ is even and, for every
$1\le k\le n$,
\[
  \sum_{i=1}^{k}d_i\;\le\;k(k-1)+\sum_{i=k+1}^{n}d_i .
\]
\end{theorem}

\begin{remark}
By Theorem~\ref{thm:main}, the sequence $(3,2,1,0)$ is not a degree sequence;
consistently, it violates the inequality of Theorem~\ref{thm:EG} at $k=2$,
since $5>2+1$.
\end{remark}

\begin{thebibliography}{9}
\bibitem{Bol} B.~Bollob\'as, \emph{Modern Graph Theory}, Graduate Texts in
Mathematics 184, Springer, 1998.
\bibitem{EG} P.~Erd\H{o}s and T.~Gallai, \emph{Gr\'afok el\H{o}\'irt fok\'u
pontokkal}, Mat. Lapok \textbf{11} (1960), 264--274.
\end{thebibliography}

\end{document}
